\scanpage{95}

\uline{Morphisme de type fini} \deleted{[\unsure{prés. finie|quasi-fini}]} d'ens schématiques : tel qu'il existe une partition de $X$ en ensembles loc. fermés constructibles affines $X_i$, telle que les $f\mid X_i$ soient des morphismes de type fini \deleted{[\unsure{prés. finie|quasi-finis}]}

% slanted
\margin{est-il indiqué de définir la présentation finie ?}

Un tel morphisme de \unsure{'ens|l'ens|un ens.} schématique est isom \unsure{:|à} un morphisme de type fini de schémas affines, Considéré comme morphisme d'ens schématiques). \uline{Transitivité}.

Ens schématique sur un autre, de type fini sur un autre. Tous sont sur $\mathbb{Z}$ !

Ens schématique défini par une partie constructible d'un ens. schématique, son immersion canonique, décomposition de $X$ en somme $X=E\amalg(X-E)$, critère de factorisabilité par une partie constructible ; c'est \uline{ensembliste} !

Produit\unsure{s} fibrés [et $\underset{\leftarrow}{\lim}$ finies] dans la catégorie des ens schématiques. \uline{Existe}, et le foncteur \deleted{\unsure{préschémas}} $(\text{présch }\unsure{\text{q. compact}|\text{qu compact}}\ \text{q. sép})\longrightarrow(\text{ens schém})$ commute aux produits fibrés.

Si on a $X$ et $Y$ sur $Z$, si $X',Y'$ sont des parties constructibles de $X,Y$ alors $X'\times_Z Y'=\mathrm{pr}_1^{-1}(X')\cap\mathrm{pr}_2^{-1}(Y')$ \ldots.
% archive sheet number: 94
